\section*{Hoofdstuk \ref{ch:g7:volcanoes} --- Vulkanen}

\begin{solution}{exo:g7:volcanoes:1}
(a) \omterm{def:g7:volcanoes:vent}{krater}; (b) \omterm{def:g7:volcanoes:magma}{magma}; (c) \omterm{def:g7:volcanoes:vent}{eruptieopening}; (d) \omterm{def:g4:volcanoes-earthquakes-first-look:eruption}{lava}; (e) \omterm{def:g7:volcanoes:magma-chamber}{magmakamer}.
\end{solution}

\begin{solution}{exo:g7:volcanoes:2}
(a) Effusief. (b) Explosief. (c) Effusief: een lavafontein werpt gloeiende
druppels uit die terugvallen en wegstromen; het \omterm{def:g7:volcanoes:magma}{magma} wordt niet tot as
geblazen. (d) Explosief: het is een \omterm{def:g7:volcanoes:pyroclastic-flow}{pyroclastische stroom}.
\end{solution}

\begin{solution}{exo:g7:volcanoes:3}
(a) Waar. (b) Niet waar: het meest voorkomende vulkanische gas is
waterdamp; koolstofdioxide komt op de tweede plaats. (c) Niet waar: een
\omterm{def:g7:volcanoes:pyroclastic-flow}{pyroclastische stroom} is een lawine van heet gas, as en stukken \omterm{def:g3:stones-pebbles-sand:rock}{gesteente}.
(d) Niet waar: ze staan midden in de Pacifische \omterm{def:g7:plate-tectonics:plate}{plaat}, meer dan 3200~km van
de dichtstbijzijnde \omterm{def:g7:plate-tectonics:plate-boundary}{plaatgrens}, boven een \omterm{def:g7:volcanoes:hotspot}{hotspot}.
\end{solution}

\begin{solution}{exo:g7:volcanoes:4}
(a) Ongeveer 1000 tot 1200~graden. (b) Ongeveer 600 tot 800~graden. (c) Bij
de \omterm{def:g7:volcanoes:vent}{eruptieopening}, waar de \omterm{def:g4:volcanoes-earthquakes-first-look:eruption}{lava} geel is (1000 tot 1200~graden) en niet
oranje (800 tot 1000~graden): de \omterm{def:g4:volcanoes-earthquakes-first-look:eruption}{lava} koelt af terwijl ze van de
\omterm{def:g7:volcanoes:vent}{eruptieopening} wegstroomt.
\end{solution}

\begin{solution}{exo:g7:volcanoes:5}
(a) $838 + 221 = 1059$ \omterm{def:g4:volcanoes-earthquakes-first-look:volcano}{vulkanen}. (b) $1059 \div 1214 \approx 0.872$:
ongeveer \textbf{87~\%}. (c) $69 + 18 + 12 + 0.5 = 99.5$, ongeveer 100~\%:
het kleine verschil komt doordat elk percentage is afgerond.
\end{solution}

\begin{solution}{exo:g7:volcanoes:6}
(a) $564 \div 660 \approx 0.855$: ongeveer \textbf{85~\%}. (b)
$660 - 564 - 62 - 33 = 1$: één steile kegel. (c) Steile kegels worden
opgebouwd door stroperige \omterm{def:g4:volcanoes-earthquakes-first-look:eruption}{lava} en as, het kenmerk van \omterm{def:g7:volcanoes:eruption-types}{explosieve
uitbarstingen}: langs de Ring van Vuur, waar een \omterm{def:g7:plate-tectonics:plate}{plaat} onder een andere
zakt, verwachten we veel \omterm{def:g7:volcanoes:eruption-types}{explosieve uitbarstingen}, zoals die van Mount
St.~Helens en de Pinatubo.
\end{solution}

\begin{solution}{exo:g7:volcanoes:7}
(a) $1170 - 1140 = 30$~graden. (b) $30 \div 10 = 3$~graden per
kilometer. (c) In de tunnel raakt de \omterm{def:g4:volcanoes-earthquakes-first-look:eruption}{lava} de lucht niet: het verharde dak
erboven houdt de warmte vast, zoals het deksel van een pan. In de open
lucht verliest het oppervlak meteen warmte aan de lucht.
\end{solution}

\begin{solution}{exo:g7:volcanoes:8}
(a) Kauai, met minstens $5.5 - 0.7 = 4.8$~miljoen jaar. (b) Kauai: \omterm{def:g4:clouds-rain-snow:precipitation}{regen} en
golven slijten het al veel langer uit. (c) Big Island, waar nog steeds
nieuwe \omterm{def:g4:volcanoes-earthquakes-first-look:eruption}{lava} naar buiten komt. (d) Naar het noordwesten: Kauai, dat als
eerste van de \omterm{def:g7:volcanoes:hotspot}{hotspot} werd weggevoerd, ligt er nu ten noordwesten van.
\end{solution}

\begin{solution}{exo:g7:volcanoes:9}
(a) Bijna alle \omterm{def:g4:volcanoes-earthquakes-first-look:volcano}{vulkanen} staan op \omterm{def:g7:plate-tectonics:plate-boundary}{plaatgrenzen} (87~\% in de lijst):
\omterm{def:g4:volcanoes-earthquakes-first-look:volcano}{vulkanen} die zo ver van elke \omterm{def:g7:plate-tectonics:plate-boundary}{plaatgrens} liggen, zijn een uitzondering.
(b) Diep in de \omterm{def:g6:inside-the-earth:layers}{mantel} onder Hawaï blijft een \omterm{def:g7:volcanoes:hotspot}{hotspot} miljoenen jaren heet.
Die doet het \omterm{def:g3:stones-pebbles-sand:rock}{gesteente} erboven gedeeltelijk smelten. Het \omterm{def:g7:volcanoes:magma}{magma}, lichter dan
het \omterm{def:g3:stones-pebbles-sand:rock}{vaste gesteente} eromheen, stijgt op door de \omterm{def:g7:plate-tectonics:plate}{plaat} en barst uit.
Terwijl de \omterm{def:g7:plate-tectonics:plate}{plaat} beweegt, groeien er nieuwe \omterm{def:g4:volcanoes-earthquakes-first-look:volcano}{vulkanen} boven de \omterm{def:g7:volcanoes:hotspot}{hotspot} en
worden de oude meegevoerd. (c) Oranje stippen (``binnen een \omterm{def:g7:plate-tectonics:plate}{plaat}'').
\end{solution}

\begin{solution}{exo:g7:volcanoes:10}
Het eerste \omterm{def:g7:volcanoes:magma}{magma} had op weg naar boven het grootste deel van zijn gas
verloren, zoals een fles bruisend water die open of stil heeft gestaan:
toen het aan het oppervlak kwam, was er niets meer over om het uiteen te
blazen, en het sijpelde naar buiten. Het \omterm{def:g7:volcanoes:magma}{magma} van 12 en 15~juni hield
zijn gas nog vast, zoals de geschudde fles: toen de druk dicht bij het
oppervlak daalde, vormde het gas in één keer bellen, die opzwollen en het
stroperige \omterm{def:g7:volcanoes:magma}{magma} tot as verscheurden. Het eerste \omterm{def:g7:volcanoes:magma}{magma} is het \omterm{def:g7:volcanoes:magma}{magma} dat
zijn gas had verloren.
\end{solution}

\begin{solution}{exo:g7:volcanoes:11}
(a) $10 \div 80 = 0.125$~uur, en $0.125 \times 60 = 7.5$~minuten. (b)
$10 \div 130 \approx 0.077$~uur, ongeveer $0.077 \times 60 \approx
4.6$~minuten, dus ongeveer 5~minuten. (c) Het gezin heeft 20~minuten nodig,
maar de stroom komt na 5 tot 8~minuten aan: als de stroom te zien is, is
het te laat. Het gezin kan alleen veilig zijn als het vertrekt zodra de
wetenschappers de waarschuwingssignalen zien, vóór de \omterm{def:g4:volcanoes-earthquakes-first-look:eruption}{uitbarsting}.
\end{solution}

\begin{solution}{exo:g7:volcanoes:12}
(a) $111 \times 360 = 39\,960$: ongeveer \textbf{40\,000~km}. (b)
$40\,000 \div 15 \approx 2667$: ongeveer 2\,670~km per dag. (c)
$2667 \div 24 \approx 111$: ongeveer \textbf{110~km per uur}. (d) Door de
wind: de fijne as heeft geen eigen snelheid meer zodra de \omterm{def:g4:volcanoes-earthquakes-first-look:eruption}{uitbarsting} haar
omhoog heeft geworpen; ze zweeft hoog in de lucht en de wind voert haar
mee, terwijl ze langzaam daalt.
\end{solution}

\begin{solution}{pb:g7:volcanoes:1}
\textbf{1.} Van 16 tot 31~maart: $31 - 16 = 15$~dagen; dan 30~dagen in
april en 18~dagen in mei: $15 + 30 + 18 = 63$~dagen.

\textbf{2.} $10\,000 \div 62 \approx 161$: gemiddeld meer dan 160 kleine
\omterm{def:g4:volcanoes-earthquakes-first-look:earthquake}{aardbevingen} per dag.

\textbf{3.} Het opzwellen liet zien dat \omterm{def:g7:volcanoes:magma}{magma} hoog in de \omterm{def:g1:landscapes-around-us:landform}{berg} was
opgestegen en het \omterm{def:g3:stones-pebbles-sand:rock}{gesteente} naar buiten duwde. Het maakte de noordflank
ook steil en instabiel, klaar om weg te glijden.

\textbf{4.} Een blauw driehoekje staat waar een \omterm{def:g7:plate-tectonics:plate}{plaat} onder een andere
zakt. We verwachten stroperig, gasrijk \omterm{def:g7:volcanoes:magma}{magma}, \omterm{def:g7:volcanoes:eruption-types}{explosieve uitbarstingen} en
een steile kegel (een stratovulkaan) --- precies wat de eerste foto laat
zien.

\textbf{5.} $240 \div 60 = 4$~kilometer per minuut.

\textbf{6.} De opgezwollen flank hield het \omterm{def:g7:volcanoes:magma}{magma} onder druk binnen, zoals
de dop van een fles. Toen die weggleed, daalde de druk in één keer; het gas
dat in het \omterm{def:g7:volcanoes:magma}{magma} was opgelost, vormde in één keer bellen en barstte naar
buiten door de open flank --- zijwaarts, waar het deksel verdwenen was.

\textbf{7.} $27 \div 480 \approx 0.056$~uur, en $0.056 \times 60 \approx
3.4$~minuten. Omdat de explosie minstens zo snel ging, deed ze er
hoogstens ongeveer 3.4~minuten over.

\textbf{8.} $8 \div 130 \approx 0.062$~uur, en $0.062 \times 60 \approx
3.7$: ongeveer \textbf{4~minuten}.

\textbf{9.} $24 \div 8.8 \approx 2.7$: de kolom steeg ongeveer 2.7 keer
hoger dan de Mount Everest.

\textbf{10.} Effusief: de stroperige \omterm{def:g4:volcanoes-earthquakes-first-look:eruption}{lava} kwam naar buiten zonder te
exploderen en stapelde zich op boven de \omterm{def:g7:volcanoes:vent}{eruptieopening}. Net als de fles
die stil had gestaan, had dit \omterm{def:g7:volcanoes:magma}{magma} het grootste deel van zijn gas al
verloren.

\textbf{11.} Voor de \omterm{def:g4:volcanoes-earthquakes-first-look:eruption}{uitbarsting} was de \omterm{def:g1:landscapes-around-us:landform}{berg} een gladde, puntige kegel
bedekt met \omterm{def:g4:clouds-rain-snow:precipitation}{sneeuw}, boven groene bossen. Daarna zijn de top en de
noordflank verdwenen: een enorme \omterm{def:g7:volcanoes:vent}{krater} die naar ons toe openstaat, heeft
hun plaats ingenomen, en de hellingen en het \omterm{def:g1:landscapes-around-us:landform}{dal} ervoor zijn grijs van as
en brokken \omterm{def:g3:stones-pebbles-sand:rock}{gesteente}.

\textbf{12.} $2950 - 2549 = 401$: de \omterm{def:g1:landscapes-around-us:landform}{berg} verloor ongeveer
\textbf{401~m} hoogte, ongeveer 400~meter, zoals in jaar~4 al werd verteld.
\end{solution}
